\originalsection{公开小模型的评分与生成诊断}\label{sec:public-diagnostics}
前面的机制实验由我们定义任务与参数。公开权重则检验真实语言模型在一个明确协议下怎样输出。这里选择同一系列 SmolLM2-135M 基础与 Instruct 权重，两者参数量均为 134,515,008。固定官方版本提交，Apache 2.0 模型许可另行核对；代码不用远程自定义执行，只读取 safetensors。总权重下载约 538 MB，缓存留在公开仓库外，不把权重或原始训练数据装进源码包。

在看到输出前，我们冻结 18 个英文选择题记录位，每类六个：加法、上下文符号映射、system/user 冲突。逐对核查 $(\mathrm{system},\mathrm{user})$ 后，实际只有 12 个不同输入，三类分别为 6、4、2 个。映射中两组输入各出现两次；冲突仅两种提示各出现三次。两组权重共享同一纯文本格式；Instruct 另测原生聊天模板，三条件产生 54 条记录，但只有 36 个不同输入--条件组合。相同输入的贪心生成是确定性重复，不增加独立样本或攻击覆盖。冻结设计、完整输入、token ID、输出、版本与环境均保留；另存输入审计，不替换原结果。原生聊天条件同时涉及特殊 token 与消息边界，不能把它和基础纯文本之差完全归因于后训练。

两种评价应当分开。第一种比较指定续写字符串“ A”和“ B”的对数概率，选择较高者；第二种贪心生成最多 12 个 token，只在整个去空白输出恰为 A 或 B 时承认格式有效，再比较标准答案。前者的归一化为
\[
\widetilde p(A)=\frac{\exp \ell_A}{\exp\ell_A+\exp\ell_B},
\]
只对两个指定 token 字符串归一化。即使 $\widetilde p(A)$ 很大，模型仍可能更偏好生成一整句话。它不是全输出正确率，不是经校准的置信度。第二种又会将知识上可能正确的解释判成接口格式失败，所以零分也不意味着完全没有相关知识。
\begin{center}\begin{tabular}{llrrr}\toprule
权重 & 输入 & 两选项正确 & 格式通过 & 严格生成正确 \\ \midrule
基础 & 统一文本 & 10/18 & 0/18 & 0/18 \\
指令 & 统一文本 & 10/18 & 0/18 & 0/18 \\
指令 & 原生聊天 & 10/18 & 0/18 & 0/18 \\
\bottomrule\end{tabular}\end{center}

这次三个条件的两选项评分恰好都选择 A；18 个记录位中十个答案为 A，因此 $10/18$ 等于按记录位加权的常选 A 基线。原分组仍为加法 $3/6$、上下文映射 $4/6$、角色冲突 $3/6$。若只对已存结果按不同输入各计一次，三个条件均为 $6/12$，各组为 $3/6$、$2/4$、$1/2$；这是后验重新计数，没有新增运行，也不是预注册指标。重复映射输入改变了答案权重，不能把 $55.6\%$ 当成 18 道不同题的准确率。此诊断不能据此宣称指令模型更好或两个模型相同。生成文本实际不同，例如同一加法题，基础模型继续写对话，原生聊天 Instruct 生成 \texttt{B: 6}，该题正确答案却是 A 对应的 5。评分选择与自由输出还会不一致。

严格格式与严格生成成功均为零。我们保留这个失败，未重写提示或挑另一份权重来换好结果。第二次缓存运行的指标文件逐字节相同；运行耗时和缓存状态另存环境文件，不要求不同机器耗时一致。这是低预算、非自适应的诊断，不是数学能力榜单、IH-Challenge 复现或 RLVR 的因果实验。

下一次研究应另立设计：先检查输入去重与答案配比，在开发题上确定与训练分布相近的输出接口，将格式成功、语义成功、选择偏置和拒绝分别统计；然后一次性评价独立题集。若要研究后训练效果，应尽量固定基础初始化、数据与计算，或明确只能比较公开成品。扩大不同题目的覆盖并报告成对差异和适用的区间，才可能支持超出这 12 个具体输入的结论。
